\begin{table}[htbp]
\centering
\caption{Diagnostic of imputed covariates}
\label{tab:diagnostico_imputacao}
\begin{tabular}{lcccc}
\toprule
Grade-subject & Complete N & Imputed N & Complete baseline & Imputed baseline \\
\midrule
9th LP & 13.992 & 15.487 & 230.07 & 229,75 \\
9th Mat & 13.999 & 15.495 & 246,56 & 245,87 \\
HS3 LP & 1.032 & 27.842 & 266,68 & 265,49 \\
HS3 Mat & 1.033 & 27.857 & 272,48 & 269,72 \\
\bottomrule
\end{tabular}
\vspace{0.2cm}
\begin{flushleft}
\footnotesize Note: complete observations have all raw covariates used in the no-imputation DR-DiD. Baseline differences across columns indicate that the no-imputation diagnostic also changes sample composition. Baseline is mean proficiency on the SARESP scale.
\end{flushleft}
\end{table}



